\documentclass[11pt]{article}

\usepackage[a4paper,margin=1in]{geometry}
\usepackage{amsmath,amssymb,amsthm,mathtools}
\usepackage{enumitem}

\newtheorem{theorem}{Theorem}
\newtheorem{lemma}{Lemma}
\newtheorem{proposition}{Proposition}
\newtheorem{corollary}{Corollary}
\newtheorem{remark}{Remark}

\title{Maximizing a Ratio Under the Constraint $0<K<C$}
\author{}


\begin{document}

\maketitle

\section{Problem setup}


We consider the general expression
\[
    \lambda(I)(\alpha)=\frac{\sum_{i=1}^{m} C^{max}_{i}(S^{WORST})+(1-2(1-\alpha))\sum_{j=1}^{n}p_{j}^{2}}{\sum_{i=1}^{m} C^{max}_{i}(S^{OPT})+(1-2(1-\alpha))\sum_{j=1}^{n}p_{j}^{2}}
\]

And let \[
    C:=\sum_{i=1}^{m} C^{max}_{i}(S^{OPT}), \qquad K:=\sum_{j=1}^{n}p_{j}^{2}.
\]

Thus, 
\[
    \lambda(I)(\alpha)=\frac{\sum_{i=1}^{m} C^{max}_{i}(S^{WORST})+(1-2(1-\alpha))K}{C+(1-2(1-\alpha))K}.
\]
We know, from Schwiegelshohn(insert reference here), that for \(\alpha \ge
\frac12, \space \text{we have:} \space \lambda(I)(\alpha) = \frac{\sum_{i=1}^{m} C^{max}_{i}(S^{WORST})}{C}\le \frac43 \).
Which implies for any feasible instance 

\[
    \sum_{i=1}^{m} C^{max}_{i}(S^{WORST})
    \leq
    \frac43 C.
\]

If we consider the function
 \[
    R(C^{max}_{i}(S^{WORST}))
    :=
    \frac{C^{max}_{i}(S^{WORST}) + \bigl(1-2(1-\alpha)\bigr)K}
         {C + \bigl(1-2(1-\alpha)\bigr)K},
\]
then $R$ is monotonically increasing in $C^{max}_{i}(S^{WORST})$. Therefore, to consider the maximum, we may consider from now only the case: 
\[
    R(C,K)(\alpha)
    :=
    \frac{\frac43 C + \bigl(1-2(1-\alpha)\bigr)K}
         {C + \bigl(1-2(1-\alpha)\bigr)K}.
\]
Furthermore, we know that for any feasible instance, we have $C>0$ and $0 \le K \le C$. And further know, 
that the case $K=C$ implies that there are exactly $m$ jobs which gives $\lambda(I)(\alpha) = 1$. The case 
$K=0$ implies that all jobs have zero processing time, which is trivial.
Therefore, we can restrict our attention to $0<K <C$.


\section{Reformulating the objective}


We have

\begin{align*}
    R(C,K,\alpha)
    &=
    \frac{\frac43 C+(1-2(1-\alpha))K}{C+(1-2(1-\alpha))K} \\
    &=
    \frac{\frac43 C+(2\alpha-1)K}{C+(2\alpha-1)K}\\
    &=
    1+
    \frac{\frac13 C}{C+(2\alpha-1)K}.
\end{align*}
The denominator is positive, so the expression is well-defined.


Because $C>0$, define
\[
    x:=\frac{K}{C}.
\]
The constraint $0<K<C$ is equivalent to
\[
    0<x<1.
\]
Using $K=xC$, we obtain
\begin{align*}
    R(C,K,\alpha)
    &=
    1+
    \frac{\frac13 C}{C-(1-2\alpha) K} \\
    &=
    1+
    \frac{\frac13 C}{C-(1-2\alpha)xC} \\
    &=
    1+
    \frac{1}{3(1-(1-2\alpha) x)}.
\end{align*}
Thus $R(C,K,\alpha)$ depends on $C$ and $K$ only through the ratio
\[
    x=\frac KC.
\]

\begin{proposition}[Monotonicity in $K/C$]
Fix $0<\alpha<1/2$. Then the function
\[
    x\mapsto 1+\frac{1}{3(1-(1-2\alpha)x)}
\]
is strictly increasing on $(0,1)$.
\end{proposition}

\begin{proof}
Let $\beta := 1-2\alpha$. Since $0<\alpha<1/2$, we have $0<\beta<1$. Define
\[
    g(x):=1+\frac{1}{3(1-\beta x)}.
\]
Then
\[
    g'(x)
    =
    \frac{\beta}{9(1-\beta x)^2}.
\]
Since $\beta>0$ and $1-\beta x>0$ for $0<x<1$, we get
\[
    g'(x)>0.
\]
Therefore $g$ is strictly increasing on $(0,1)$.
\end{proof}

\section{Supremum for $0<\alpha<1/2$}

\begin{theorem}[Supremum of the ratio]
Let
\[
    C>0, \qquad 0<K<C, \qquad 0<\alpha<\frac12.
\]
Then
\[
    R(C,K,\alpha)<1+\frac{1}{6\alpha}
\]
\end{theorem}

\begin{proof}
From the monotonicity of $R$ in $x=K/C$, we know that the supreum is reached only for $x \rightarrow 1$. \\
So, 

\begin{align*}
    R(C,K,\alpha)
    &=
    1+
    \frac{1}{3(1-(1-2\alpha)x)}\\
    &<
    1+
    \frac{1}{3(1-(1-2\alpha))}.\\
    &=
    1+\frac{1}{6\alpha}
\end{align*}

\end{proof}

\end{document}
